\section{Vector Fields}
A \textcolor{Bittersweet}{\textbf{vector field}} on $\mathbb{R}^{n}$ is a
function $\mathbf{F}:\mathbb{R}^{n}\to\mathbb{R}^{n}$. A
\textcolor{Bittersweet}{gradient field} is the field of gradients
\[\nabla f\left(x,y,z\right)=\left<f_x\left( x,y,z \right), f_y\left( x,y,z
\right), f_z\left( x,y,z \right)  \right>.\]
where $f$ is the potential function of the gradient field.
The \textcolor{Bittersweet}{line integral} along segment $C$ is given by
\[\int_C f ds = \int_{{b}}^{{a}} {f\left(\mathbf{r}\left( t \right) \right)} \left| \mathbf{v}\left( t \right)  \right| \: d{t},\quad \text{for some smooth $\mathbf{r}$}\] 
The line integral is dependent on the path $C$, but not on the motion
$\mathbf{r},\mathbf{v}$. It could be interpreted as the \textcolor{Blue}{area
under the path}.
\paragraph{Properties of Vector Fields}
Let $\mathbf{F}$ be a vector field on an open connected domain $D$.\\
If $\forall A,B\in D$, 
$\int_{{A}}^{{B}} {\mathbf{F}\cdot} d{\mathbf{r}} {}=f(B)-f(A)$
is the same over all paths from $A$ to $B$, the integral is
\textcolor{Bittersweet}{path independent} in $D$ and $\mathbf{F}$ is
\textcolor{Bittersweet}{conservative} on $D$. If $\mathbf{F}=\nabla f$ for some
scalar function $f:\mathbb{R}^{n}\to \mathbb{R}$, $f$ is a potential function
for $\mathbf{F}$.
\begin{align*}
  \mathbf{F}\text{ is a gradient field}\iff& \mathbf{F}\text{ is a conservative
field}\\
    \iff& \oint_C \mathbf{F}\cdot d{\mathbf{r}=0}\\
    \iff& curl\:{\mathbf{F}}=\nabla \times \mathbf{F}=\left<0,0,0 \right> 
\end{align*}

\paragraph{Applications}
The \textcolor{Bittersweet}{work done} by a force $\mathbf{F}$ over a smooth
curve $\mathbf{r}\left( t \right)$ is
\[W=\int_{{t=a}}^{{t=b}} {\mathbf{F\cdot T}} \: d{s} {}=\mathbf{F}\left( \mathbf{r}\left( t \right)  \right)\cdot \mathbf{r}'\left( t \right)dt.\]
The work integral is independent of the motion $\mathbf{r}$, but is path
dependent.\\
The \textcolor{Bittersweet}{flow} along a curve for some smooth $\mathbf{r}$ in
a continuous velocity field $\mathbf{F}$ is
\[Flow=\int_{{a}}^{{b}} {\mathbf{F\cdot T}} \: d{s} {}.\] 
If the curve is a closed loop, the flow integral is the
\textcolor{Bittersweet}{circulation about the curve}. Mathematically, it is
identical to the work done.

\paragraph{Green's Theorem}
A \textcolor{Bittersweet}{closed curve} is one whose starting point $=$ end
point. A \textcolor{Bittersweet}{simple curve} is one that does not cross
itself. A \textcolor{Bittersweet}{directed/oriented} closed curve is one with a
direction. An oriented simple closed curve moving in an
\textcolor{Blue}{anti-clockwise direction} has a
\textcolor{Bittersweet}{positive }orientation, and negative orientation
otherwise.\\
For a positive oriented closed curve $C$, $R$ its enclosed region, and
$\mathbf{F}\left( x,y \right) =\left<M(x,y), N(x,y) \right> $ a differentiable
vector field on the $xy$-plane
\[\oint_C M\:dx+N\:dy=\iint_R \frac{\partial N}{\partial x} -\frac{\partial M}{\partial y} \:dx\:dy.\] 
\[\oint_C M\:dy-N\:dx=\iint_R \frac{\partial M}{\partial x} +\frac{\partial N}{\partial y} \:dx\:dy.\] 
For vector fields undefined at the origin, we consider a curve $C_h$ clockwise
around the origin of radius $h$. 

\paragraph{Surface Integrals}
The area of surface $S$ $f\left( x,y,z \right) =c$ over a closed and bounded
plane region $R$ is, where $\mathbf{p}$ is a the unit vector normal to $R$,
\[Surface\ Area=\iint_R \frac{\left| \nabla f \right| }{\left| \nabla f\cdot
\mathbf{p} \right| }dA.\]
For a continuous $g$ defined at points of $S$, the surface integral is
\[\iint_S g d\sigma=\iint_R g\left( x,y,z \right) \frac{\left| \nabla f \right| }{\left|
\nabla f\cdot \mathbf{p} \right| }dA.\] 

Intuitively, an \textcolor{Bittersweet}{orientation} of a surface is defining a
unit normal vector $\mathbf{n}\left( p \right) $ for every point $p$ on a
surface, such that the vector changes \textcolor{Blue}{continuously} wrt.  $p$.
A surface with an orientation is orientable. A connected surface has either $0$
or $2$ orientations, and closed surfaces are orientable.\\
The \textcolor{Bittersweet}{flux} of a vector field $\mathbf{F}$ across an
oriented surface $S$ in the direction of $\mathbf{n}$ is
\[Flux=\iint_S\mathbf{F}\cdot d\mathbf{S}=\iint_S \mathbf{F}\left( x,y,z \right) \cdot \mathbf{n}\:d\sigma.\]
Intuitively, it is the ``amount of flowing'' across a surface. For a surface $S$ defined
by $g(x,y,z)=0$, 
\[Flux=\iint_R \mathbf{F}\cdot \frac{\nabla g}{\left| \nabla g\cdot  \left<0,
0, 1 \right>  \right| }dx\:dy\] 

\paragraph{Stokes' Theorem}
We define the curl the ``amount of spinning'' in a vector field. We use the
right-hand rule for the direction of the curl.
\[curl\: \mathbf{F}=\nabla\times \mathbf{F}=\left<\frac{\partial P}{\partial y} -\frac{\partial N}{\partial z} ,\frac{\partial M}{\partial z} -\frac{\partial P}{\partial x} ,\frac{\partial N}{\partial x} -\frac{\partial M}{\partial y}  \right>.\] 
By Green's theorem, $\oint_C \mathbf{C}\cdot d\mathbf{r}=\iint_R curl\: \mathbf{F}\cdot \mathbf{k}dA$
The circulation (flow integral) a vector field $\mathbf{F}=\left<M,N,P \right>
$ around the boundary $C$ of an oriented surface $S$ with $\mathbf{n}$ in the
counter-clockwise direction
\[\oint_C \mathbf{F}\cdot d\mathbf{r}=\iint_S\nabla\times \mathbf{F\cdot n}d\sigma.\] 

\paragraph{Divergence Theorem}
We define the divergence the ``amount of expansion/production'' in a vector field.
\[div\: \mathbf{F}=\nabla\cdot\mathbf{F}=\frac{\partial M}{\partial x}
+\frac{\partial N}{\partial y} +\frac{\partial P}{\partial z} .\] 
The field is incompressible if $div = 0$ at all points. The theorem states
\[ \iint_S\mathbf{F\cdot n}\:d\sigma=\iiint_D div\:\mathbf{F}dV\]
